% © 2026 Ing. David Jaroš — CC BY-NC-ND 4.0
\documentclass[12pt,a4paper]{article}
\usepackage[a4paper,margin=2.5cm]{geometry}
\usepackage{amsmath,amssymb,amsthm,mathtools}
\usepackage[most]{tcolorbox}
\newtheorem{definition}{Definition}[section]
\newtheorem{theorem}[definition]{Theorem}
\newtheorem{corollary}[definition]{Corollary}
\newtheorem{proposition}[definition]{Proposition}
\theoremstyle{remark}\newtheorem{remark}[definition]{Remark}
\newcommand{\B}{\mathbb B}
\title{GAP-10I Pairing Audit:\\No-New-Field Reduction and the Concurrent-Vector Obstruction}
\author{Ing. David Jaro\v{s}}
\date{19 July 2026}
\begin{document}
\maketitle

\begin{abstract}
The generic two-sided UBT derivative
$D_\mu X=\partial_\mu X+A_\mu X-XB_\mu$ appears to introduce two new
connection fields.  This note proves that this is not necessary in the pure
Lorentz branch: preservation of the intrinsic Lorentz real slice and its
metric reduces the pair uniquely, modulo a central term that cancels from the
bimodule derivative, to the single existing spin connection,
$D_\mu X=\partial_\mu X+\Omega_\mu X+X\Omega_\mu^\ddagger$.
However, the same reduction exposes a decisive obstruction.  If
$E_\mu=D_\mu\Theta/\sqrt{\mathcal N_0}$ is a nondegenerate torsion-free tetrad,
then a Lorentz-real projection of $\Theta$ defines a concurrent spacetime
vector $V$ satisfying $\nabla_\mu V^\nu=\delta_\mu{}^\nu$.  Hence the metric
admits a proper homothety and $R^\rho{}_{\sigma\mu\nu}V^\sigma=0$.
The non-flat Schwarzschild vacuum exterior with nonzero mass, and generic
torsion-free Einstein geometries, do not obey this restriction.  Thus the
no-new-field Lorentz pairing is rigorously closed as a kinematic reduction and
its torsion-free generated-tetrad branch is closed as a no-go.  This statement
is not a no-go for the same pairing with composite torsion: the companion note
\texttt{gap\_10i\_torsionful\_local\_representer.tex} constructs an explicit
local right inverse for every smooth tetrad.  A relative bimodule action is
therefore optional for local kinematics, though it remains a possible
torsion-free completion if derived without propagating degrees of freedom.
\end{abstract}

\section{Algebraic conventions}
Let $\B\simeq\operatorname{Mat}(2,\mathbb C)$ and let $\ddagger$ denote
Hermitian conjugation (complex conjugation followed by quaternion
conjugation).  The classical Lorentz slice is
\begin{equation}
 W_L=\{X\in\B\mid X^\ddagger=-X\}
 =\{i x^0\mathbf1+x^k\mathbf e_k\mid x^a\in\mathbb R\}.
\end{equation}
Equivalently, $W_L$ is the fixed set of
$\mathcal JX=-X^\ddagger$.

Consider the most general regular bimodule connection on $\B$,
\begin{equation}
 D_\mu X=\partial_\mu X+A_\mu X-XB_\mu.
 \label{eq:general-pair}
\end{equation}

\section{Exact reduction of the pair}
\begin{theorem}[Lorentz-slice preserving pairs]
At a fixed spacetime point, Eq.~\eqref{eq:general-pair} maps $W_L$ into
$W_L$ for every $X\in W_L$ if and only if
\begin{equation}
 \boxed{B_\mu=-A_\mu^\ddagger+c_\mu\mathbf1,\qquad c_\mu\in\mathbb R.}
 \label{eq:slice-pair}
\end{equation}
\end{theorem}

\begin{proof}
For $X^\ddagger=-X$, slice preservation is equivalent to
$(D_\mu X)^\ddagger=-D_\mu X$.  The connection terms give
\begin{equation}
 (A_\mu+B_\mu^\ddagger)X
 =X(A_\mu^\ddagger+B_\mu).
 \label{eq:P-condition}
\end{equation}
Put $P_\mu=A_\mu+B_\mu^\ddagger$.  Taking $X=i\mathbf1$ in
Eq.~\eqref{eq:P-condition} gives $P_\mu=P_\mu^\ddagger$.  Taking
$X=\mathbf e_k$ then gives $[P_\mu,\mathbf e_k]=0$ for all three quaternion
units.  Their commutant in $\operatorname{Mat}(2,\mathbb C)$ is the center,
so $P_\mu=c_\mu\mathbf1$; Hermiticity makes $c_\mu$ real.  This is
Eq.~\eqref{eq:slice-pair}.  The converse follows by direct substitution.
\end{proof}

\begin{theorem}[Metric-compatible reduction to one spin connection]
Assume in addition that the induced transport preserves the Lorentz quadratic
form.  Write
\begin{equation}
 A_\mu=\Omega_\mu+\alpha_\mu\mathbf1,
 \qquad \operatorname{tr}\Omega_\mu=0.
\end{equation}
Then metric compatibility requires
$c_\mu=2\operatorname{Re}\alpha_\mu$, and Eq.~\eqref{eq:general-pair}
reduces exactly to
\begin{equation}
 \boxed{D_\mu X=\partial_\mu X+\Omega_\mu X+X\Omega_\mu^\ddagger.}
 \label{eq:lorentz-pair}
\end{equation}
Equivalently one may choose the representative
\begin{equation}
 \boxed{A_\mu=\Omega_\mu,\qquad B_\mu=-\Omega_\mu^\ddagger.}
 \label{eq:A-B-Omega}
\end{equation}
The common central one-form is invisible because it cancels between left and
right multiplication.
\end{theorem}

\begin{proof}
Using Eq.~\eqref{eq:slice-pair}, the connection part becomes
\begin{equation}
 A_\mu X-XB_\mu
 =\Omega_\mu X+X\Omega_\mu^\ddagger
 +(2\operatorname{Re}\alpha_\mu-c_\mu)X.
\end{equation}
The final real scalar term is a local dilation of every Lorentz vector and
therefore rescales the Lorentz bilinear form.  Metric compatibility sets its
coefficient to zero.  Then
$B_\mu=-\Omega_\mu^\ddagger+\alpha_\mu\mathbf1$, and the same central
$\alpha_\mu\mathbf1$ occurs on both sides of Eq.~\eqref{eq:general-pair}, so
it cancels identically.  The traceless spin lift fixes the representative
Eq.~\eqref{eq:A-B-Omega}.
\end{proof}

\begin{corollary}[No additional Lorentz connection fields]
In the pure gravitational Lorentz branch, $A_\mu$ and $B_\mu$ are not two
independent fields.  They are the left and right actions of the one spin
connection
\begin{equation}
 \Omega_\mu=\frac14\omega_{\mu ab}\Sigma^{ab}.
\end{equation}
For specified tetrad and torsion,
$\Omega=\Omega(e,T)$ is already fixed.  In the torsion-free branch it is the
Levi--Civita spin connection, and in Cartesian inertial Minkowski gauge
$A_\mu=B_\mu=\Omega_\mu=0$.
\end{corollary}

\section{The obstruction exposed by the reduction}
\begin{theorem}[Concurrent-vector theorem]
Let the torsion-free metric-compatible derivative be
Eq.~\eqref{eq:lorentz-pair}.  Suppose
\begin{equation}
 D_\mu\Theta=sE_\mu,\qquad s:=\sqrt{\mathcal N_0},
 \qquad E_\mu\in W_L,
 \label{eq:tetrad-generation}
\end{equation}
and $E_\mu$ is a nondegenerate tetrad.  Then there exists a real spacetime
vector field $V$ satisfying
\begin{equation}
 \boxed{\nabla_\mu V^\nu=\delta_\mu{}^\nu.}
 \label{eq:concurrent}
\end{equation}
Consequently
\begin{equation}
 \mathcal L_Vg_{\mu\nu}=2g_{\mu\nu},
 \qquad
 R^\rho{}_{\sigma\mu\nu}V^\sigma=0.
 \label{eq:homothety-curvature}
\end{equation}
\end{theorem}

\begin{proof}
Equation~\eqref{eq:lorentz-pair} commutes with
$\mathcal JX=-X^\ddagger$.  Therefore the Lorentz-real projection
\begin{equation}
 X:=\frac12(\Theta+\mathcal J\Theta)\in W_L
\end{equation}
obeys $D_\mu X=sE_\mu$, because $\mathcal JE_\mu=E_\mu$.  Write
$X=X^a\mathbf u_a$ and $E_\mu=e_\mu{}^a\mathbf u_a$ in a Lorentz basis and
define
\begin{equation}
 V^\nu:=s^{-1}e_a{}^\nu X^a.
\end{equation}
Using the inverse tetrad postulate and $D_\mu X^a=s e_\mu{}^a$ gives
\begin{equation}
 \nabla_\mu V^\nu
 =s^{-1}e_a{}^\nu D_\mu X^a
 =e_a{}^\nu e_\mu{}^a
 =\delta_\mu{}^\nu.
\end{equation}
Symmetrization yields the homothety equation.  Antisymmetrizing two covariant
derivatives of $V$ yields the curvature-kernel identity.
\end{proof}

\begin{proposition}[Torsion-free Schwarzschild-exterior exclusion]
The non-flat Schwarzschild vacuum exterior with mass $M\ne0$ admits no vector
satisfying Eq.~\eqref{eq:concurrent}; indeed it admits no proper homothety
$\mathcal L_Vg=2g$.
\end{proposition}

\begin{proof}
Average a putative homothetic vector over the rotational isometry group.  The
averaged vector remains homothetic and has the form
$V=T(t,r)\partial_t+R(t,r)\partial_r$.  For
\begin{equation}
 ds^2=-f(r)dt^2+f(r)^{-1}dr^2+r^2d\Omega^2,
 \qquad f(r)=1-\frac{2M}{r},
\end{equation}
the angular component of $\mathcal L_Vg=2g$ gives $R=r$.  The $rr$ component
then gives
\begin{equation}
 -\frac{r f'(r)}{f(r)^2}+\frac{2}{f(r)}=\frac{2}{f(r)},
\end{equation}
so $f'(r)=0$, contradicting $f'(r)=2M/r^2\ne0$.
\end{proof}

\begin{remark}[Constants, auxiliary fields, and the torsionful escape]
A central constant or central one-form cancels from the bimodule derivative
and cannot change Eq.~\eqref{eq:concurrent}.  A potential may make an
additional field auxiliary or fix its vacuum value, but a constant parameter
alone cannot provide a spacetime-dependent torsion-free relative action.
However, the obstruction also disappears when the same one-connection pair
uses a nonzero metric-compatible composite contortion.  That local construction
is given in \texttt{gap\_10i\_torsionful\_local\_representer.tex}.  Neither
route is dynamically selected by the present kinematic theorem.
\end{remark}

\section{Defensible gap status}
\begin{tcolorbox}[colback=green!6!white,colframe=green!55!black,title=New exact subclosures]
\textbf{GAP-10I-PAIR-KIN: CLOSED [L1].}
Lorentz-slice and metric compatibility reduce the apparent fields
$(A_\mu,B_\mu)$ to the single existing spin connection, modulo a central term
that cancels identically.  In flat inertial gauge both vanish.

\medskip
\textbf{GAP-10I-PAIR-GR: CLOSED AS A TORSION-FREE NO-GO [L1].}
That pure Lorentz pairing with $K=0$ makes a nondegenerate generated tetrad
imply a concurrent vector and a proper homothety.  It therefore excludes the
non-flat Schwarzschild vacuum exterior with $M\ne0$ and generic torsion-free
Einstein geometries.

\medskip
\textbf{GAP-10I-TORSION-LOCAL: CLOSED LOCALLY [L1] IN A COMPANION THEOREM.}
The same one-connection pairing with explicit composite metric-compatible
contortion has a local single-$\Theta$ representer for every smooth Lorentzian
tetrad on a non-null Gaussian patch.

\medskip
\textbf{GAP-10I-2S: NOT REQUIRED FOR LOCAL KINEMATICS.}
A relative bimodule connection remains an optional torsion-free completion and,
if used, must be derived as composite or auxiliary rather than propagating.

\medskip
\textbf{GAP-10I-CURVED: LOCAL KINEMATICS CLOSED; DYNAMICS/GLOBAL PART
NARROWED.}
Canonical action selection, physical torsion constraints, regularity at
$V^2=0$, horizons/caustics, and global continuation remain open.

\medskip
\textbf{GAP-10D: NOT CLOSED BY THIS REDUCTION.}
The result identifies the torsion-free obstruction but does not choose between
the composite-contortion branch and an optional torsion-free relative action.
The Palatini/Lovelock endpoint remains conditional until canonical UBT derives
the selected branch, its algebraic elimination or spin current, and the
induced metric variation.
\end{tcolorbox}

\end{document}
